%% manuscript.tex — Lancet Regional Health – Americas format
%% Overleaf-ready standalone version
%% Compilar com: pdflatex → pdflatex (bibliografia embutida)
%% ===================================================================
\documentclass[review,number]{elsarticle}

\usepackage[utf8]{inputenc}
\usepackage[T1]{fontenc}
\usepackage{amsmath,amssymb}
\usepackage{graphicx}
\usepackage{booktabs}
\usepackage{multirow}
\usepackage{xcolor}
\usepackage{lineno}
\usepackage{hyperref}
\usepackage[margin=2.5cm]{geometry}

\biboptions{super,comma}   % Superscript Vancouver citations (Lancet style)

\linenumbers  % Obrigatório para submissão

%% Metadados do journal
\journal{The Lancet Regional Health -- Americas}

\begin{document}

\begin{frontmatter}

%% ── Título ──────────────────────────────────────────────────────────
\title{Cold, not heat: delayed temperature extremes drive respiratory
       hospitalisations in subtropical Southern Brazil}

%% ── Autores ─────────────────────────────────────────────────────────
\author[ufpr]{Felipe Baglioli\corref{cor1}}
\ead{felipe.baglioli@ufpr.br}
\cortext[cor1]{Corresponding author}

\author[ufpr]{Lucas Rover}

\address[ufpr]{Federal University of Paraná (UFPR), Curitiba, PR, Brazil}

%% ── Abstract estruturado (formato Lancet, ≤300 palavras) ───────────
\begin{abstract}

\textbf{Background}\quad
Climate--health policy in the Southern Hemisphere has prioritised heat
alerts, yet the respiratory burden attributable to cold extremes in
subtropical cities remains largely unquantified. This evidence gap
leaves millions without targeted early-warning systems for cold spells.

\textbf{Methods}\quad
We analysed daily minimum temperature (INMET, 1961--2024) and
respiratory hospital admissions (DATASUS SIH/AIH, ICD-10 J00--J99,
2022--2024; $N$=42\,716) in Curitiba, Brazil ($\sim$1${\cdot}$8~million
inhabitants) using a methodological triangulation design. Distributed
lag non-linear models (DLNM) with quasi-Poisson regression provided
inferential evidence; K-means clustering assessed whether the same
delayed patterns emerge without parametric assumptions; and
gradient-boosted trees (XGBoost) with SHAP feature importance provided
predictive convergence evidence. The MMT was identified
by grid search. Sensitivity analyses varied lag windows, seasonal
degrees of freedom, and excluded 2022 for residual COVID-19 effects.

\textbf{Findings}\quad
Over 1096~study days, 42\,716~respiratory hospitalisations (51\% male;
38\% aged $\geq$65~years) and 3715~in-hospital deaths were recorded.
The DLNM identified MMT=17${\cdot}$8\,°C. Cold extremes at the 1st
percentile of $T_\mathrm{min}$ (5${\cdot}$2\,°C) yielded cumulative
RR=1${\cdot}$42 [95\% CI: 1${\cdot}$09--1${\cdot}$84] over
0--21~days --- corresponding to approximately 14 excess respiratory
admissions per extreme cold day; at the 5th percentile
(8${\cdot}$5\,°C), RR=1${\cdot}$22 [1${\cdot}$02--1${\cdot}$44].
Effects were delayed, peaking at lags 3--10~days. In contrast, heat
effects were non-significant (P90 RR=1${\cdot}$06
[0${\cdot}$98--1${\cdot}$14]). Results were robust across all
sensitivity analyses ($\hat{\phi}$=1${\cdot}$32, Ljung--Box
$p$=0${\cdot}$14). Two independent methods converged on the same
finding: the ``Clean Cold'' cluster exhibited delayed correlations
with hospitalisations at lags 3--10~days without imposing any
parametric lag structure, and SHAP analysis ranked $T_\mathrm{min}$
as the dominant environmental predictor.

\textbf{Interpretation}\quad
In this subtropical metropolis, cold extremes --- not heat --- are the
dominant temperature-related driver of respiratory hospitalisations,
with a delayed effect structure that opens a 3--10-day window for
preventive intervention. The convergence of three methodologically
independent approaches --- inferential (DLNM),
structural-exploratory (clustering), and predictive (ML/SHAP) ---
strengthens this conclusion. Current early-warning systems in
subtropical South America focus almost exclusively on heat alerts;
integrating cold-spell forecasts could reduce a substantial and
preventable respiratory burden.

\textbf{Funding}\quad
None.

\end{abstract}

%% ── Keywords ────────────────────────────────────────────────────────
\begin{keyword}
temperature extremes \sep cold spells \sep hospital admissions \sep
distributed lag non-linear model \sep methodological triangulation \sep
climate--health \sep respiratory morbidity \sep subtropical \sep Curitiba
\end{keyword}

\end{frontmatter}

%% ═══════════════════════════════════════════════════════════════════
%% RESEARCH IN CONTEXT (formato obrigatório Lancet)
%% ═══════════════════════════════════════════════════════════════════

\section*{Research in Context}

\textbf{Evidence before this study}

We searched PubMed and Scopus for studies published between
January~2010 and December~2025 using the terms ``temperature
extremes'', ``hospital admissions'', ``distributed lag'', and
``Southern Hemisphere'' or ``Latin America''. The landmark
384-location study by Gasparrini et~al.\ (2015) established U-shaped
temperature--mortality curves globally, yet included no subtropical
South American city. Subsequent multi-city analyses have reinforced
that cold-attributable mortality exceeds heat-attributable mortality
in most climates, but morbidity evidence from the Southern Hemisphere
remains scarce. No study has applied formal DLNM methods to quantify
the non-linear, delayed effects of cold extremes on respiratory
hospital admissions in this region.

\textbf{Added value of this study}

To our knowledge, this is the first DLNM analysis of temperature
extremes and respiratory hospitalisations in a subtropical South
American city. We show that cold extremes at the 1st percentile of
$T_\mathrm{min}$ (5${\cdot}$2\,°C) are associated with a 42\%
increase in respiratory hospitalisations (RR=1${\cdot}$42
[1${\cdot}$09--1${\cdot}$84]), while heat effects are
non-significant. The delayed peak at 3--10~days reveals a clinically
actionable lag structure not previously documented with inferential
methods in this region. These findings challenge the heat-centric
framing that dominates climate--health policy in South America.
Crucially, we employ methodological triangulation: unsupervised
clustering identifies the same delayed cold--hospitalisation pattern
without parametric assumptions, and SHAP-based feature importance
from a gradient-boosted tree model independently ranks
$T_\mathrm{min}$ as the dominant environmental predictor. This
three-method convergence strengthens the evidence beyond any single
analytical framework.

\textbf{Implications of all the available evidence}

Early-warning systems in subtropical Southern Hemisphere cities should
integrate cold-spell forecasts alongside heat alerts. The 3--10-day
delay between cold exposure and hospitalisation peak provides a
concrete window for preventive action --- including targeted outreach
to elderly and respiratory-vulnerable populations, temporary shelter
activation, and primary-care surge planning. Given that current
Brazilian climate--health surveillance (VIGIAR-SUS) lacks a
cold-specific trigger, our findings, corroborated by three independent analytical
approaches, identify a remediable gap in public health preparedness.

%% ═══════════════════════════════════════════════════════════════════
%% INTRODUÇÃO
%% ═══════════════════════════════════════════════════════════════════
\section{Introduction}

Global evidence increasingly shows that cold-attributable mortality
exceeds heat-attributable mortality in most
climates,\cite{Gasparrini2015} yet climate--health policy remains
overwhelmingly focused on heat alerts and heatwave preparedness.
This asymmetry is particularly pronounced in subtropical cities of
the Southern Hemisphere, where winter temperatures can fall below
5\,°C but early-warning systems rarely include cold-spell
triggers.\cite{Bell2024,McMichael2006}

The pathophysiology of cold-related respiratory morbidity is well
established. Low ambient temperatures suppress mucociliary clearance,
impair alveolar macrophage function, and increase indoor crowding ---
conditions that favour respiratory pathogen
transmission.\cite{Linares2025} Cold also triggers systemic
vasoconstriction, elevates blood viscosity, and raises cardiovascular
strain, compounding the respiratory burden in vulnerable
populations.\cite{Ebi2021} Critically, unlike heat-related
morbidity, which manifests acutely within 0--3~days, cold-related
effects exhibit delayed lags of 3--14~days, complicating their
detection in routine surveillance and their attribution in ecological
studies.\cite{Cheng2024}

Despite these known mechanisms, formal inferential evidence on the
temperature--morbidity relationship remains scarce for subtropical
South America. The landmark 384-location study by Gasparrini
et~al.\cite{Gasparrini2015} did not include any subtropical Southern
Brazilian city. Multi-city studies in Brazil have focused on air
pollution and mortality,\cite{Requia2024} while temperature--morbidity
analyses have relied on correlation-based approaches that cannot
disentangle delayed effects from shared seasonality. As a result,
a fundamental question remains unanswered for this region: does cold
or heat impose the greater respiratory hospitalisation burden, and
over what time horizon?

We address this gap using distributed lag non-linear models (DLNM)
applied to 42\,716~respiratory hospital admissions and daily minimum
temperature records in Curitiba, the largest subtropical city in
southern Brazil ($\sim$1${\cdot}$8~million inhabitants, $\sim$930~m
elevation, Cfb climate). Curitiba's climate --- with winter minima
below 5\,°C and summer maxima approaching 30\,°C --- provides a
natural experiment for contrasting cold and heat effects within a
single population. Our primary objective is to quantify the
non-linear, delayed exposure--response relationship between
temperature extremes and respiratory admissions, with a pre-specified
hypothesis that cold effects would dominate in this setting.
We employ a methodological triangulation design: the DLNM provides
inferential evidence; unsupervised clustering tests whether the same
delayed patterns emerge from the data structure without parametric
assumptions; and gradient-boosted tree models with SHAP
feature-importance analysis offer an independent, predictive line
of convergence.

%% ═══════════════════════════════════════════════════════════════════
%% MÉTODOS
%% ═══════════════════════════════════════════════════════════════════
\section{Methods}

\subsection{Study design and setting}

We conducted an ecological time-series study of daily environmental
exposures and respiratory hospital admissions in Curitiba, Paraná,
Brazil (population~$\sim$1${\cdot}$8~million) from
1~January~2022 to 31~December~2024 (1096~days). Curitiba is the
largest city in southern Brazil with a Köppen Cfb (oceanic subtropical)
climate, winter minima below 5\,°C, and summer maxima approaching
30\,°C --- making it representative of the broader subtropical belt
of southern South America. The study-period $T_\mathrm{min}$
distribution (mean 14${\cdot}$8\,°C, P10=9${\cdot}$8\,°C) was
comparable to the 1961--2024 INMET climatology (mean 13${\cdot}$3\,°C,
P10=7${\cdot}$8\,°C; eFigure~9), indicating that the 2022--2024
period did not experience a regime shift in cold extremes.

\subsection{Data sources}

Individual-level hospital admission records were obtained from the
Brazilian Unified Health System (SIH/SUS--DATASUS) for 2022--2024
($N$=42\,716 respiratory admissions, ICD-10 J00--J99). The primary
outcome was the daily count of respiratory admissions. Demographics:
51\% male, 38\% aged $\geq$65~years. In-hospital mortality was
8${\cdot}$7\% ($n$=3715).

Daily minimum ($T_\mathrm{min}$), mean, and maximum temperatures,
relative humidity, and atmospheric pressure were obtained from INMET
automatic station~A807. For ETCCDI percentile calculations, the full
record (1961--2024) was used; combined analyses were restricted to
2022--2024.

PM$_{2{\cdot}5}$ concentrations were measured using a PurpleAir
Flex-II sensor (Sensor ID~90875) at daily resolution. Raw readings
were corrected using the EPA nationwide correction
equation:\cite{Barkjohn2021}
%
\begin{equation}
  \mathrm{PM}_{2{\cdot}5}^{\mathrm{corr}} =
    \text{0·524} \times \mathrm{PM}_{2{\cdot}5}^{\mathrm{sensor}}
    - \text{0·0862} \times \mathrm{RH} + \text{5·75}
  \label{eq:epa}
\end{equation}
%
This correction has been validated against Federal Equivalent Method
monitors\cite{Holder2020} and is widely adopted for PurpleAir
data.\cite{Magi2020,Barkjohn2022} An inter-channel quality control
protocol excluded days when inter-channel differences exceeded
5~µg/m$^3$ and 70\% relative difference simultaneously. After QC,
697~days (63${\cdot}$6\%) had valid EPA-corrected PM$_{2{\cdot}5}$.

\subsection{Definitions and missing data}

Extreme temperature days were defined using ETCCDI monthly percentiles
(TX90p for extremely hot, TN10p for extremely cold) over the 1961--2024
baseline. Heatwaves: $\geq$3~consecutive extreme hot days; cold spells:
$\geq$3~consecutive extreme cold days. Missing data rates:
$T_\mathrm{min}$~2${\cdot}$6\%, PM$_{2{\cdot}5}$~36${\cdot}$4\%,
hospitalisations~0${\cdot}$2\%. Primary analyses used complete cases;
sensitivity analyses with imputation produced similar results
(Supplement).

\subsection{Clustering analysis (structural-exploratory evidence)}

To assess whether the delayed cold--hospitalisation pattern emerges
from the multivariate data structure without parametric assumptions,
K-means clustering was applied to six standardised variables
($T_\mathrm{max}$, $T_\mathrm{med}$, $T_\mathrm{min}$,
PM$_{2{\cdot}5}$, relative humidity, atmospheric pressure) on
complete-case days, with optimal $K$ selected by silhouette score.
Lagged cross-correlation functions (lags 0--28~days) were computed
between cluster membership and daily hospitalisations. PCA was
performed on standardised PM$_{2{\cdot}5}$, $T_\mathrm{min}$,
and hospitalisations ($N$=683). This analysis does not impose a
distributed lag function, spline basis, or risk-ratio framework, and
thus constitutes an independent test of the lag structure identified
by the DLNM (eTables~2--5, eFigures~1--4). Time-series were also
decomposed using STL, and ADF tests assessed stationarity.

\subsection{Distributed lag non-linear model}
\label{sec:dlnm}

The primary analysis used distributed lag non-linear models
(DLNM),\cite{Gasparrini2010,Gasparrini2011} which simultaneously
capture the non-linear shape of the exposure--response curve and the
delayed lag structure --- a critical requirement given the known
3--14-day latency of cold-related respiratory effects:
%
\begin{equation}
  \log\!\big(\mathrm{E}[Y_t]\big) = \alpha
    + \mathbf{s}\bigl(x_t, \mathrm{lag}; \boldsymbol{\eta}\bigr)
    + \sum_{d=1}^{6} \beta_d \,\mathrm{DOW}_{d,t}
    + \gamma \,\mathrm{Holiday}_t
    + \mathrm{ns}(\mathrm{time}_t,\, 7\,\mathrm{df/year})
  \label{eq:dlnm}
\end{equation}
%
where $Y_t$ is the daily count of respiratory admissions (quasi-Poisson),
$\mathbf{s}(\cdot)$ is the cross-basis function (natural cubic splines:
exposure df=4, lag df=5), DOW denotes day-of-week indicators, and
$\mathrm{ns}(\mathrm{time}, 21\,\mathrm{df})$ captures trend and
seasonality. Maximum lag was 21~days. The minimum morbidity temperature
(MMT) was identified by grid search, following
Gasparrini et al.\cite{Gasparrini2015} Confidence intervals were
computed via the delta method. Sensitivity analyses varied the maximum
lag (14, 21, 28~days), seasonal df (6, 7, 8~df/year), and excluded
2022 to assess residual COVID-19 effects.

As a third, methodologically independent line of evidence, machine
learning models (XGBoost,\cite{Chen2016} multilayer
perceptron,\cite{Goodfellow2016} extreme learning
machine\cite{Huang2006}) were trained to forecast daily
hospitalisations at lags 0--7~days, with SHAP
values\cite{Lundberg2017} quantifying feature importance. Because
gradient-boosted trees use recursive partitioning rather than spline
bases and optimise predictive accuracy rather than estimating relative
risks, agreement between SHAP rankings and DLNM findings constitutes
methodological convergence. Full details are reported in the
Supplement (eMethods, eTable~8, eFigure~8).

\paragraph{Ethics and reporting}
This study used exclusively aggregated, publicly available data from
DATASUS, INMET, and PurpleAir. Under Brazilian CNS~510/2016,
research using public aggregated data is exempt from ethics review.
Reporting follows the STROBE guidelines for observational
studies.\cite{VonElm2007}

%% ═══════════════════════════════════════════════════════════════════
%% RESULTADOS
%% ═══════════════════════════════════════════════════════════════════
\section{Results}

\subsection{Descriptive statistics}

Over 1096~study days, we recorded 42\,716~respiratory hospital
admissions and 3715~in-hospital deaths (case fatality 8${\cdot}$7\%),
representing a mean of 32${\cdot}$3~admissions per day
(SD~11${\cdot}$7, range 1--73; Table~\ref{tab:descriptive}). The
study population was predominantly male (51\%) and elderly (38\%
aged $\geq$65~years). EPA-corrected PM$_{2{\cdot}5}$ was generally
low (median~4${\cdot}$12~µg/m$^3$), with only 4${\cdot}$7\% of days
exceeding the WHO 24-hour guideline (15~µg/m$^3$). Using ETCCDI
monthly percentiles, we identified 215~extremely hot days
(20${\cdot}$1\%) and 39~extremely cold days (3${\cdot}$7\%),
aggregating into 25~heatwave events but only 4~cold-spell events
(eTable~1) --- underscoring the episodic and concentrated nature of
cold exposure in this climate.

\begin{table}[ht]
\centering
\caption{Descriptive statistics for all study variables, Curitiba,
2022--2024 ($N$=1096~days).}
\label{tab:descriptive}
\begin{tabular}{lrrrrrrrr}
\toprule
Variable & $N$ & Mean & SD & Median & Q25 & Q75 & Min & Max \\
\midrule
$T_\mathrm{max}$ (°C) & 1070 & 24${\cdot}$53 & 4${\cdot}$75 & 24${\cdot}$80 & 21${\cdot}$50 & 28${\cdot}$10 & 10${\cdot}$20 & 35${\cdot}$40 \\
$T_\mathrm{med}$ (°C) & 1042 & 18${\cdot}$74 & 3${\cdot}$63 & 18${\cdot}$90 & 16${\cdot}$30 & 21${\cdot}$50 &  7${\cdot}$80 & 27${\cdot}$30 \\
$T_\mathrm{min}$ (°C) & 1068 & 14${\cdot}$75 & 3${\cdot}$75 & 15${\cdot}$20 & 12${\cdot}$00 & 17${\cdot}$70 &  1${\cdot}$30 & 23${\cdot}$40 \\
PM$_{2{\cdot}5}$ (µg/m$^3$) &  697 &  5${\cdot}$84 & 7${\cdot}$40 &  4${\cdot}$12 &  2${\cdot}$43 &  6${\cdot}$55 &  0${\cdot}$00 & 65${\cdot}$19 \\
Rel.\ humidity (\%)     &  859 & 79${\cdot}$98 &  9${\cdot}$52 & 80${\cdot}$90 & 74${\cdot}$15 & 86${\cdot}$55 & 43${\cdot}$00 & 98${\cdot}$40 \\
Hosp.\ (n/day)   & 1094 & 32${\cdot}$29 & 11${\cdot}$74 & 32${\cdot}$00 & 23${\cdot}$00 & 41${\cdot}$00 & 1${\cdot}$00 & 73${\cdot}$00 \\
\bottomrule
\end{tabular}
\end{table}

\subsection{Clustering: structural-exploratory evidence}

K-means clustering on six climate--pollution variables ($K$=3,
silhouette=0${\cdot}$357) identified three environmental regimes:
Clean Cold ($n$=248), Polluted Heat ($n$=12), and Clean Heat
($n$=298; eTables~2--3, eFigures~1--4). The silhouette score is
moderate, which is typical of continuous environmental data where
clusters represent gradients rather than discrete
categories;\cite{Rousseeuw1987} a validation analysis at $K$=4
yielded stronger between-cluster differences in hospitalisations
(Kruskal--Wallis $H$=16${\cdot}$69, $p<$0${\cdot}$001; eTable~3).
Crucially, the Clean Cold cluster exhibited delayed positive
correlations with hospitalisations peaking at lags 3--10~days
--- the same lag window subsequently identified by the DLNM ---
without imposing any distributional lag function or spline basis.
PCA corroborated this pattern: PC1 (44${\cdot}$5\% of variance)
loaded positively on PM$_{2{\cdot}5}$ and hospitalisations and
negatively on $T_\mathrm{min}$ (eTable~4), indicating that cold,
polluted days are structurally associated with higher admissions.
Extended lag analysis confirmed peak
$T_\mathrm{min}$--hospitalisation correlation at lag~9
($r$=$-$0${\cdot}$320, $p<$0${\cdot}$001; eTable~5). This
convergence demonstrates that the delayed cold effect is a
structural feature of the data, not an artefact of the DLNM
parametric specification.

\subsection{Distributed lag non-linear models}
\label{sec:dlnm_results}

\paragraph{Temperature exposure--response}
The quasi-Poisson DLNM for $T_\mathrm{min}$ ($N$=1068,
$\hat{\phi}$=1${\cdot}$32) identified a minimum morbidity temperature
(MMT) of 17${\cdot}$8\,°C, closely aligned with the 75th percentile
of the local temperature distribution
(Figure~\ref{fig:dlnm_temp_overall}). Below this threshold, the
cumulative exposure--response curve rose steeply and monotonically:
at the 1st percentile of $T_\mathrm{min}$ (5${\cdot}$2\,°C),
RR=1${\cdot}$42 [1${\cdot}$09--1${\cdot}$84] --- equivalent to a
42\% excess in respiratory admissions; at the 5th percentile
(8${\cdot}$5\,°C), RR=1${\cdot}$22 [1${\cdot}$02--1${\cdot}$44]; at
the 10th percentile (9${\cdot}$8\,°C), RR=1${\cdot}$17
[1${\cdot}$00--1${\cdot}$37] (Table~\ref{tab:dlnm_temp_rr}). In
stark contrast, heat effects were non-significant across all upper
percentiles (P90 RR=1${\cdot}$06 [0${\cdot}$98--1${\cdot}$14];
P99 RR=1${\cdot}$29 [0${\cdot}$98--1${\cdot}$68]). The lag
structure at cold (P10) peaked at 3--10~days
(Figure~\ref{fig:dlnm_temp_lag}a), consistent with the known latency
of cold-induced respiratory pathology and with the cross-correlation
analysis.

\paragraph{Sensitivity and diagnostics}
The cold effect was robust to all pre-specified sensitivity analyses
(Table~\ref{tab:sensitivity}). Cumulative RR at P10 ranged from
1${\cdot}$05 (lag=14~days, capturing only the early phase of the
delayed response) to 1${\cdot}$23 (lag=28~days). Notably, excluding
2022 --- the year with the greatest residual COVID-19 influence ---
strengthened the cold effect to RR=1${\cdot}$37
[1${\cdot}$11--1${\cdot}$70], suggesting that pandemic-related
disruptions attenuated the temperature--morbidity signal in the full
period. ADF tests confirmed non-stationarity in raw hospitalisations
($p$=0${\cdot}$164) but stationarity in STL residuals
($p$=0${\cdot}$0003; eFigure~7), validating the detrending approach.
Model diagnostics showed no residual autocorrelation
(Durbin--Watson=2${\cdot}$00, Ljung--Box $p$=0${\cdot}$14;
eTable~7).

\begin{table}[ht]
\centering
\caption{Cumulative relative risk (RR) at key percentiles of
$T_\mathrm{min}$, centred at MMT=17${\cdot}$8\,°C. Quasi-Poisson DLNM,
lag 0--21~days, $N$=1068.}
\label{tab:dlnm_temp_rr}
\begin{tabular}{lrrr}
\toprule
Percentile & $T_\mathrm{min}$ (°C) & RR & 95\% CI \\
\midrule
P1  &  5${\cdot}$2 & 1${\cdot}$418 & 1${\cdot}$091--1${\cdot}$842 \\
P5  &  8${\cdot}$5 & 1${\cdot}$215 & 1${\cdot}$023--1${\cdot}$442 \\
P10 &  9${\cdot}$8 & 1${\cdot}$167 & 0${\cdot}$998--1${\cdot}$365 \\
P25 & 12${\cdot}$0 & 1${\cdot}$133 & 0${\cdot}$979--1${\cdot}$312 \\
P75 & 17${\cdot}$7 & 1${\cdot}$001 & 0${\cdot}$996--1${\cdot}$005 \\
P90 & 19${\cdot}$4 & 1${\cdot}$055 & 0${\cdot}$978--1${\cdot}$138 \\
P95 & 20${\cdot}$0 & 1${\cdot}$102 & 0${\cdot}$979--1${\cdot}$240 \\
P99 & 21${\cdot}$4 & 1${\cdot}$285 & 0${\cdot}$983--1${\cdot}$679 \\
\bottomrule
\end{tabular}
\end{table}

\begin{table}[ht]
\centering
\caption{Sensitivity analyses for the temperature DLNM. Cumulative RR
at the 10th percentile of $T_\mathrm{min}$ (9${\cdot}$8\,°C).}
\label{tab:sensitivity}
\begin{tabular}{lrrrr}
\toprule
Scenario & RR & 95\% CI & $\hat{\phi}$ & $N$ \\
\midrule
Main (lag=21, 7\,df/yr) & 1${\cdot}$17 & 1${\cdot}$00--1${\cdot}$37 & 1${\cdot}$32 & 1068 \\
Lag=14                  & 1${\cdot}$05 & 0${\cdot}$92--1${\cdot}$18 & 1${\cdot}$31 & 1068 \\
Lag=28                  & 1${\cdot}$23 & 1${\cdot}$01--1${\cdot}$51 & 1${\cdot}$32 & 1068 \\
Seasonal 6\,df/yr       & 1${\cdot}$17 & 1${\cdot}$00--1${\cdot}$36 & 1${\cdot}$33 & 1068 \\
Seasonal 8\,df/yr       & 1${\cdot}$21 & 1${\cdot}$03--1${\cdot}$42 & 1${\cdot}$31 & 1068 \\
Exclude 2022            & 1${\cdot}$37 & 1${\cdot}$11--1${\cdot}$70 & 1${\cdot}$26 &  718 \\
\bottomrule
\end{tabular}
\end{table}

\begin{figure}[ht]
\centering
\includegraphics[width=0.85\textwidth]{figures/fig_dlnm_temp_overall.pdf}
\caption{Cumulative exposure--response curve for $T_\mathrm{min}$
and respiratory hospitalisations (DLNM, lag 0--21~days). The curve
rises steeply below the MMT (17${\cdot}$8\,°C, red dashed line),
indicating a cold-dominant pattern. Heat effects above the MMT are
non-significant. Shaded area: 95\% CI.}
\label{fig:dlnm_temp_overall}
\end{figure}

\begin{figure}[ht]
\centering
\includegraphics[width=\textwidth]{figures/fig_dlnm_temp_lag.pdf}
\caption{Lag-specific RR for (a) cold extreme (P10, 9${\cdot}$8\,°C)
and (b) heat extreme (P90, 19${\cdot}$4\,°C), centred at
MMT=17${\cdot}$8\,°C.}
\label{fig:dlnm_temp_lag}
\end{figure}

\subsection{Machine learning: predictive convergence evidence}

XGBoost achieved the best forecasting performance
(MAPE=19${\cdot}$4\%, RMSE=8${\cdot}$37 at lag~0; eTable~8).
SHAP feature-importance analysis ranked $T_\mathrm{min}$ as the
dominant environmental predictor of respiratory hospitalisations
(eFigure~8). This is noteworthy because gradient-boosted trees use
recursive partitioning --- a fundamentally different approach from
the DLNM's cross-basis splines --- yet $T_\mathrm{min}$ emerges as
the top-ranked predictor, constituting methodological convergence.
Together, the three approaches form a triangulation:
\begin{itemize}
  \item \textbf{DLNM:} causal-ecological evidence (non-linear RR with lag structure);
  \item \textbf{Clustering:} structural-exploratory evidence (delayed cold--hospitalisation pattern without parametric assumptions);
  \item \textbf{ML/SHAP:} predictive convergence evidence ($T_\mathrm{min}$ as dominant predictor via tree-based partitioning).
\end{itemize}

%% ═══════════════════════════════════════════════════════════════════
%% DISCUSSÃO
%% ═══════════════════════════════════════════════════════════════════
\section{Discussion}

\subsection{Principal findings}

This study provides the first formal DLNM-based evidence that cold
extremes --- not heat --- are the dominant temperature-related driver
of respiratory hospitalisations in a subtropical South American city.
At the 1st percentile of $T_\mathrm{min}$ (5${\cdot}$2\,°C),
cumulative RR reached 1${\cdot}$42 [1${\cdot}$09--1${\cdot}$84] over
0--21~days, corresponding to approximately 14 excess respiratory
admissions per extreme cold day. Heat effects, by contrast, were
non-significant across all upper percentiles. The delayed lag peak
at 3--10~days identifies a clinically actionable window between cold
exposure and peak hospitalisation burden.

\subsection{Cold-dominant burden: challenging the heat-centric
paradigm}

A dominant framing in climate--health research equates rising
temperatures with rising health risk.\cite{Bell2024,McMichael2006}
Our findings complicate this narrative for subtropical settings. In
Curitiba, the asymmetry between cold and heat effects was striking:
the cold extreme (P1) yielded a statistically significant 42\%
excess, while the heat extreme (P90) produced a non-significant 6\%
excess. This pattern is consistent with the global observation by
Gasparrini et~al.\cite{Gasparrini2015} that cold-attributable
mortality exceeds heat-attributable mortality in most climates, but
extends it to respiratory morbidity in a Southern Hemisphere
subtropical city not represented in that 384-location study. Our
MMT of 17${\cdot}$8\,°C --- closely matching the local 75th
percentile --- aligns with subtropical MMTs reported in multi-country
analyses, and the magnitude of cold-attributable risk at P1 is
comparable to cold-related excess mortality in temperate European
and East Asian settings.\cite{Gasparrini2015}

The practical consequence is clear: in subtropical cities, a
heat-exclusive early-warning system will miss the largest
temperature-attributable morbidity signal.

\subsection{Pathophysiological plausibility and lag structure}

The 3--10-day lag peak is biologically coherent. Cold exposure
initiates a cascade that unfolds over days: immediate airway cooling
suppresses mucociliary clearance and reduces epithelial barrier
integrity (hours); subsequent impairment of alveolar macrophage
phagocytosis and secretory IgA production increases susceptibility
to viral and bacterial infection (1--3~days); clinical illness and
hospitalisation then follow after an incubation and progression
period (3--10~days).\cite{Linares2025,Ebi2021} This temporal sequence
--- unlike the rapid onset of heat-related presentations --- explains
why cold-attributable morbidity is systematically underdetected by
surveillance systems that operate on a 0--48-hour event window.

Additionally, cold weather increases indoor crowding and reduces
ventilation, amplifying respiratory pathogen transmission at the
population level.\cite{Cheng2024} These behavioural mechanisms
compound the direct physiological effects, creating a multi-pathway
exposure that is unique to cold extremes.

\subsection{Robustness of findings}

The cold effect was robust to all pre-specified sensitivity analyses.
The direction and magnitude were consistent across lag windows
(14--28~days), seasonal degrees of freedom (6--8/year), and time
periods. Notably, excluding 2022 --- the year most affected by
residual COVID-19 pandemic dynamics --- strengthened the cold effect
(RR=1${\cdot}$37 [1${\cdot}$11--1${\cdot}$70]), suggesting that
pandemic-related changes in healthcare-seeking behaviour and viral
ecology may have attenuated the temperature--morbidity signal in the
full study period. This observation implies that the true
pre-pandemic cold effect may be even larger than our main estimate.

A distinctive feature of this study is the methodological
triangulation that strengthens the cold-effect finding beyond the
DLNM alone. The K-means clustering identifies the same delayed
cold--hospitalisation pattern (lags 3--10~days) without imposing
any distributional lag function, spline basis, or risk-ratio
framework --- the pattern emerges purely from the multivariate
structure of the data. The XGBoost model, optimised for predictive
accuracy via recursive tree partitioning, independently ranks
$T_\mathrm{min}$ as the dominant environmental predictor through
SHAP analysis. These three methods differ in their assumptions
(semi-parametric splines vs.\ distance-based clustering vs.\
ensemble decision trees), optimisation criteria (quasi-likelihood
vs.\ within-cluster variance vs.\ predictive error), and outputs
(relative risks vs.\ cluster profiles vs.\ feature importance).
Their convergence on the same conclusion --- that cold stress is
the primary environmental driver of respiratory hospitalisations,
with delayed effects --- provides substantially stronger evidence
than any single analytical framework alone, and reduces the concern
that results are artefacts of a particular modelling choice.

\subsection{Policy implications}

These findings have direct implications for climate--health
preparedness in subtropical South America.\cite{Bell2024} In Brazil,
the national climate--health surveillance system (VIGIAR-SUS)
currently lacks a cold-specific trigger for respiratory early
warnings. Our results suggest that cold-spell forecasts should be
integrated alongside heat alerts, with three specific actions
enabled by the 3--10-day lag structure:

First, the delayed onset of the hospitalisation peak provides a
concrete window for \emph{preventive intervention}: targeted outreach
to elderly and respiratory-vulnerable populations can be initiated
immediately after cold-spell onset, before the morbidity wave
arrives.

Second, primary-care and emergency departments can use cold-spell
forecasts for \emph{surge planning}, anticipating increased
respiratory admissions 3--10~days after a cold event.

Third, the identification of MMT=17${\cdot}$8\,°C provides a
\emph{quantitative trigger threshold} that can be operationalised in
local early-warning systems, rather than relying on arbitrary
temperature cut-offs.

However, as a single-city study, these results require validation in
multi-city pooled analyses before informing national policy, although
the methodological triangulation (DLNM, clustering, ML/SHAP) partially
compensates by demonstrating internal convergence across independent
analytical frameworks.
Extension to other subtropical Southern Hemisphere cities following
the multi-country framework of Gasparrini
et~al.\cite{Gasparrini2015} would substantially strengthen the
evidence base.\cite{Requia2024}

\subsection{Limitations}

Several limitations warrant consideration, though none invalidate
the principal findings. First, as an ecological time-series study,
we cannot infer individual-level causal relationships; however, the
ecological design is appropriate for the study's primary objective
--- identifying population-level patterns to inform early-warning
systems --- and individual-level confounding is partially addressed
by the time-series structure.

Second, the 3-year analysis period (2022--2024) encompasses only
three seasonal cycles and four cold-spell events; however, the DLNM
estimates derive from the full continuous temperature distribution
($N$=1068~days), sensitivity analyses consistently reproduced the
cold effect (Table~\ref{tab:sensitivity}), and the convergence of
three methodologically independent approaches provides additional
assurance that the finding is not a statistical artefact of the
short series.

Third, confounding by respiratory viral circulation (influenza, RSV)
cannot be excluded; however, the fact that three analytically
independent methods --- each with different assumptions and data
subsets --- converge on the same cold-lag pattern makes it unlikely
that unmeasured viral confounding alone explains the observed
associations.

Fourth, we used total respiratory admissions (ICD-10
J00--J99) without age or cause-specific stratification;
nonetheless, the aggregate association is the clinically relevant
signal for population-level early-warning systems, and
age-stratified analyses are a priority for future work.

Fifth, temperature data from a single INMET station (A807) may not
capture microclimatic variation across the metropolitan area;
however, cold extremes are driven by synoptic-scale frontal systems
that affect the entire region more uniformly than localised urban
heat island effects.

Sixth, PM$_{2{\cdot}5}$ data from a low-cost sensor had
36${\cdot}$4\% missing values (not MCAR; eTable~10). The A-B-C
confounding diagnostic (eTable~9) proved that PM$_{2{\cdot}5}$ does
not confound the temperature association
($\Delta$RR$<$0${\cdot}$001 between Models B and C). Importantly,
the primary cold-effect finding does not depend on PM$_{2{\cdot}5}$
data: both the DLNM and ML/SHAP analyses use temperature as primary
exposure independently of PM$_{2{\cdot}5}$ completeness.

Finally, as a single-city study, results require replication;
however, the methodological triangulation partially compensates for
the absence of multi-site replication by demonstrating internal
convergence across independent analytical frameworks.

\subsection{Conclusion}

In Curitiba, cold extremes --- not heat --- are the dominant
temperature-related driver of respiratory hospitalisations, with a
42\% excess risk at the coldest days (RR=1${\cdot}$42
[1${\cdot}$09--1${\cdot}$84]) and a delayed effect peaking at
3--10~days that creates a concrete window for preventive
intervention. The convergence of three independent analytical
approaches --- inferential (DLNM), structural-exploratory
(clustering), and predictive (ML/SHAP) --- strengthens this
finding by demonstrating that the cold-lag pattern emerges from the
data regardless of modelling assumptions. These findings challenge
the heat-centric framing of climate--health surveillance in
subtropical South America and identify a specific, remediable gap:
the absence of cold-spell triggers in early-warning systems.
Integrating a temperature-threshold alert at the identified MMT
(17${\cdot}$8\,°C) into existing surveillance frameworks could
reduce a preventable respiratory burden in one of South America's
largest subtropical populations.

%% ═══════════════════════════════════════════════════════════════════
%% DECLARAÇÕES (Lancet backmatter order)
%% ═══════════════════════════════════════════════════════════════════
\section*{Contributors}
FB: Conceptualisation, Data curation, Formal analysis, Investigation,
Methodology, Writing --- original draft. LR: Methodology, Software,
Validation, Formal analysis, Writing --- review \& editing. All
authors had full access to all the data in the study and had final
responsibility for the decision to submit for publication.

\section*{Declaration of interests}
The authors declare no competing interests.

\section*{Data sharing}
All code and processed datasets are available at
\url{https://github.com/Roverlucas/cwb-temp-pollution-health}.
Raw hospitalisation data: DATASUS
(\url{https://datasus.saude.gov.br/}); meteorological data: INMET
station~A807; PM$_{2{\cdot}5}$: PurpleAir Sensor ID~90875.

\section*{Funding}
This research did not receive any specific grant from funding agencies
in the public, commercial, or not-for-profit sectors.

\section*{Acknowledgements}
The authors thank the Brazilian Institute of Meteorology (INMET),
DATASUS, and PurpleAir Inc.\ for providing open data.

%% ═══════════════════════════════════════════════════════════════════
%% REFERÊNCIAS (Vancouver, numbered by first citation)
%% ═══════════════════════════════════════════════════════════════════
\bibliographystyle{elsarticle-num}

\begin{thebibliography}{18}

%% Ordered by first citation in text (Vancouver style)

\bibitem{Gasparrini2015}
Gasparrini A, Guo Y, Hashizume M, et al.
Mortality risk attributable to high and low ambient temperature: a multicountry observational study.
\textit{Lancet}. 2015;386(9991):369--375.

\bibitem{Bell2024}
Bell ML, Gasparrini A, Benjamin GC.
Climate Change, Extreme Heat, and Health.
\textit{N Engl J Med}. 2024;390(19):1793--1801.

\bibitem{McMichael2006}
McMichael AJ, Woodruff RE, Hales S.
Climate change and human health: present and future risks.
\textit{Lancet}. 2006;367(9513):859--869.

\bibitem{Linares2025}
Linares C, et al.
How air pollution and extreme temperatures affect emergency hospital admissions due to various respiratory causes in Spain, by age group.
\textit{Int J Hyg Environ Health}. 2025;266:114570.

\bibitem{Ebi2021}
Ebi KL, et al.
Hot weather and heat extremes: health risks.
\textit{Lancet}. 2021;398(10301):698--708.

\bibitem{Cheng2024}
Cheng C, et al.
Effects of extreme temperature events on deaths and its interaction with air pollution.
\textit{Sci Total Environ}. 2024;915:170212.

\bibitem{Requia2024}
Requia WJ, et al.
Short-term air pollution exposure and mortality in Brazil: Investigating the susceptible population groups.
\textit{Environ Pollut}. 2024;340:122797.

\bibitem{Barkjohn2021}
Barkjohn KK, Gantt B, Clements AL.
Development and application of a United States-wide correction for PM2.5 data collected with the PurpleAir sensor.
\textit{Atmos Meas Tech}. 2021;14(6):4617--4637.

\bibitem{Holder2020}
Holder AL, Mebust AK, Maghran LA, et al.
Field evaluation of low-cost particulate matter sensors for measuring wildfire smoke.
\textit{Sensors}. 2020;20(17):4796.

\bibitem{Magi2020}
Magi BI, Cupini C, Francis J, et al.
Evaluation of PM2.5 measured in an urban setting using a low-cost optical particle counter and a Federal Equivalent Method Beta Attenuation Monitor.
\textit{Aerosol Sci Technol}. 2020;54(2):147--159.

\bibitem{Barkjohn2022}
Barkjohn KK, et al.
Using low-cost sensors to quantify the effects of air quality events on indoor and outdoor PM2.5 concentrations.
\textit{Sensors}. 2022;22(3):1010.

\bibitem{Gasparrini2010}
Gasparrini A, Armstrong B, Kenward MG.
Distributed lag non-linear models.
\textit{Stat Med}. 2010;29(21):2224--2234.

\bibitem{Gasparrini2011}
Gasparrini A.
Distributed lag linear and non-linear models in R: the package dlnm.
\textit{J Stat Softw}. 2011;43(8):1--20.

\bibitem{Chen2016}
Chen T, Guestrin C.
XGBoost: A scalable tree boosting system.
In: \textit{Proc 22nd ACM SIGKDD Int Conf Knowledge Discovery and Data Mining}. 2016:785--794.

\bibitem{Goodfellow2016}
Goodfellow I, Bengio Y, Courville A.
\textit{Deep Learning}.
MIT Press; 2016.

\bibitem{Huang2006}
Huang G-B, Zhu Q-Y, Siew C-K.
Extreme learning machine: Theory and applications.
\textit{Neurocomputing}. 2006;70(1--3):489--501.

\bibitem{Lundberg2017}
Lundberg SM, Lee S-I.
A unified approach to interpreting model predictions.
In: \textit{Advances in Neural Information Processing Systems}. 2017;30:4765--4774.

\bibitem{VonElm2007}
von Elm E, Altman DG, Egger M, et al.
The Strengthening the Reporting of Observational Studies in Epidemiology (STROBE) statement: guidelines for reporting observational studies.
\textit{Lancet}. 2007;370(9596):1453--1457.

\bibitem{Rousseeuw1987}
Rousseeuw PJ.
Silhouettes: A graphical aid to the interpretation and validation of cluster analysis.
\textit{J Comput Appl Math}. 1987;20:53--65.

\end{thebibliography}

\end{document}
